\section{Evaluation Methodology}
\label{sec:evaluation}

\subsection{Two non-interchangeable MQAR protocols}
\label{sec:mqar-protocols}

Multi-query associative recall (MQAR) reveals whether a model can bind keys to
values and retrieve several bindings later in a sequence. It is also easy to
mis-measure by shifting targets twice or comparing different generators.
\ALUCLU therefore ships two separately identified protocols.

\begin{table}[h]
\centering
\small
\caption{MQAR paths in the reference implementation.}
\label{tab:mqar-protocols}
\begin{tabularx}{\textwidth}{@{}p{0.19\textwidth}XX@{}}
\toprule
Property & Diagnostic next-token MQAR & Zoology-aligned MQAR \\
\midrule
Implementation &
\texttt{generate\_mqar\_batch}, \texttt{mqar\_accuracy} &
\texttt{generate\_zoology\_mqar\_batch},
\texttt{zoology\_mqar\_loss/accuracy} \\
Target meaning &
Predict the value after a query; normal LM shifting applies &
Targets are already placed at the query position; no extra shift is allowed \\
Purpose &
Compact developer learnability and regression diagnostic &
Protocol-compatible study of historical Zoology data semantics \\
Claim ceiling &
Local evidence for the exact declared distribution &
Comparable only after full dataset, grid, seed, and selection policy are run \\
\bottomrule
\end{tabularx}
\end{table}

The historical raw manifest records vocabulary 8192, training and test sizes
100,000 and 3,000, data seed zero, test-seed offset 10, power parameter 0.01,
model seed 123, weight decay 0.1, and up to 64 epochs. It records four
\((\text{length},\text{pairs},\text{batch})\) points:
\((64,4,512)\), \((128,8,512)\), \((256,16,256)\), and
\((512,64,128)\). A mismatch between a historical four-layer description and
two configuration entries is preserved in the raw manifest. The repository's
smallest runnable interpretation has a different protocol identity; it is not
silently presented as the original experiment.

\subsection{P0: correctness and small-scale decision gate}

The tests cover deterministic invariants, but the architecture should not
scale merely because those tests pass. The proposed P0 matrix in
\Cref{tab:p0} determines whether each lane earns its state and latency cost.

\begin{longtable}{@{}p{0.18\textwidth}p{0.27\textwidth}
  p{0.27\textwidth}p{0.20\textwidth}@{}}
\caption{Proposed P0 experiments. These sweeps are protocols, not current
results.}
\label{tab:p0}\\
\toprule
Experiment & Variants and sweep & Metrics & Decision rule \\
\midrule
\endfirsthead
\toprule
Experiment & Variants and sweep & Metrics & Decision rule \\
\midrule
\endhead
Partition equivalence &
All lanes and full model; token, prime-sized, random, adversarial chunks &
Maximum/mean error in output, state, and gradients &
Remain within dtype-specific preregistered tolerance \\
Causal prefix &
Replace future suffixes across seeds &
Maximum error on unchanged prefix &
No future-dependent change beyond tolerance \\
State saturation &
Streams below and far beyond retained horizon &
Logical bytes, allocated/reserved memory, post-saturation slope &
Logical persistent slope equals zero after allocation \\
Archive wrap &
Multiple ring wraps and rank sweep &
Counts, cursor, denominator minimum, reconstruction error, ledger &
No bound violation; ledger never under-reports algebraic tail \\
Cache semantics &
Increasing, decreasing, tied, outlier, duplicate, heterogeneous priorities &
Retained K/V, priorities, scan/step outputs &
Strict greater-than, batch-local policy matches oracle \\
Recall curves &
Lane subsets across distance, pair count, overwrites, capacity ratios &
Accuracy/NLL by distance and \(T/H_{\max}\) &
Complete seed distributions and curves, not selected points \\
Numerical stress &
fp32/bfloat16, long carried state, repeated compression &
NaN/Inf, norms, denominator, ledger, first failure index &
All declared paths finite or failure fully localized \\
CPU reference &
Batch 1/4/16, short and saturated streams &
Token/s, p50/p95/p99 latency, logical/resident memory &
Descriptive cost curves only \\
\bottomrule
\end{longtable}

Quality comparisons should use at least five seeds and retain all runs.
Accuracy differences should use paired bootstrap intervals or an appropriate
interval for paired proportions. Scale-up is justified only if the full model
or a simpler ablation is nondominated on a quality--state--latency frontier. A
lane that supplies no reproducible Pareto point should be removed rather than
retained for visual symmetry.

\subsection{P1: model quality and long-context transfer}

P1 begins with a matched-token 100--200M-parameter pilot. A 1B-class run is
permitted only if the pilot effect survives multiple seeds. This staged design
reduces the risk of spending scale to conceal an architectural weakness.

\begin{table}[p]
\centering
\small
\caption{Proposed model-quality matrix.}
\label{tab:quality-matrix}
\begin{tabularx}{\textwidth}{@{}p{0.19\textwidth}p{0.25\textwidth}
  p{0.27\textwidth}X@{}}
\toprule
Evaluation & Question & Sweep & Required report \\
\midrule
Held-out text NLL &
Does the memory design improve language modeling? &
Training length and controlled extrapolation &
Token-weighted NLL, perplexity, intervals, all seeds \\
Zoology MQAR &
Where does associative recall fail relative to budgets? &
Full grid, learning rates, distances, overwrites, capacity ratios &
Accuracy by point and \(T/H_{\max}\), protocol identity \\
RULER \citep{hsieh2024ruler} &
Can the model retrieve, trace, and aggregate? &
4K through 128K where supported &
Task-by-length curves, not only macro average \\
NoLiMa \citep{modarressi2025nolima} &
Does retrieval survive weak lexical overlap? &
Supported lengths, stratified overlap &
Absolute score and degradation from short-context score \\
HELMET \citep{yen2025helmet} &
Does synthetic recall transfer to application tasks? &
All supported categories and lengths &
Per-category uncertainty; synthetic kept separate \\
BABILong \citep{kuratov2024babilong} &
Does carried memory support multi-step reasoning? &
Supported long lengths and declared fine-tuning &
Task scores, training regime, confidence intervals \\
\bottomrule
\end{tabularx}
\end{table}

\subsection{Baseline contract}

The minimum rerun baseline set is:

\begin{enumerate}
  \item a causal Transformer with a mature exact-attention kernel and a
  fixed-window Transformer;
  \item Based \citep{arora2024based};
  \item pure Gated Delta and Gated DeltaNet-2
  \citep{yang2025gated,hatamizadeh2026gated};
  \item Mamba and current maintained Mamba-family controls
  \citep{gu2024mamba,dao2024transformers,lahoti2026mamba3};
  \item a reproducible residual-selected bounded-cache control, including HOLA
  if its released protocol and implementation are suitable
  \citep{cui2026hippocampus};
  \item full \ALUCLU and every principal lane subset.
\end{enumerate}

External paper numbers provide context but do not enter a comparative result
table unless rerun under the shared data, token, parameter, precision, and
hardware contract. A baseline with growing global KV state may be a valuable
quality control, but must not be labeled context-independent.

\subsection{Ablation program}

\begin{longtable}{@{}p{0.29\textwidth}p{0.35\textwidth}
  p{0.27\textwidth}@{}}
\caption{Minimum proposed ablations.}
\label{tab:ablations}\\
\toprule
Ablation & Hypothesis isolated & Primary measurements \\
\midrule
\endfirsthead
\toprule
Ablation & Hypothesis isolated & Primary measurements \\
\midrule
\endhead
Gated Delta only &
Dense recurrent baseline &
NLL, recall by distance, state, throughput \\
Local only; recurrent + local &
Exact recency and complementarity &
Near/far recall, state, latency \\
Episodic combinations &
Exact segments and compressed archives &
Recall around eviction, ledger, compression-tail latency \\
Exact-cache combinations &
Residual-selected events versus recency &
Recall by retained age, turnover, state and read cost \\
Archive off; \(K,C,A\) sweep &
Exact-live versus compressed allocation &
Capacity-boundary recall, error and bytes \\
Routing off; degree one; sketch router &
Incremental value of routing and polynomial features &
Quality, wall time, bytes, calibration \\
\(p,m,J\) sweep &
Sketch collision and compute budget &
Recall, score variance, FFT cost, stability \\
Rank \(r\) sweep &
Compression capacity &
Dense reconstruction error, ledger, p99 latency \\
Window \(W\) and cache \(C_x\) sweeps &
Exact memory allocation &
Age curves, bytes, latency, traffic \\
Uniform, fixed, learned fusion &
Token-conditional lane selection &
Quality, entropy, branch collapse, gradients \\
Residual, FIFO, recency, random, reservoir admission &
Specific value of write-residual priority &
Recall, age, diversity, turnover, runtime \\
Priority aging and diversity controls &
Cache freezing and duplicate retention &
New-information recall and unique clusters \\
fp32 versus bfloat16 state &
Precision versus traffic &
Drift, finite values, bytes and quality \\
\bottomrule
\end{longtable}

\subsection{Systems scenarios}
\label{sec:systems-evaluation}

Reference code and optimized kernels must be reported separately. The proposed
systems matrix includes interactive short prompts, interactive 16K--128K
prompts, batched throughput, saturated streaming across multiple archive
wraps, prefill scaling, and decode-only runs from saturated snapshots.
Required metrics include time to first token, p50/p95/p99 time per output
token, aggregate token/s, logical and allocated state, peak resident memory,
DRAM bytes/token, and joules/token.

Archive-boundary tokens must be tagged separately because QR/SVD can create a
tail-latency mode that the median hides. Random input generation, compilation,
autotuning, model loading, and state construction are reported as separate
regions. Warm-up, device synchronization, clock policy, and raw samples are
preserved.

\subsection{Fused-kernel correctness gate}

No fused CUDA or Triton implementation is part of the current evidence. A
future accelerator path must agree with the recurrent oracle on:

\begin{itemize}
  \item forward outputs and every carried-state tensor;
  \item input and parameter gradients;
  \item fp32, bfloat16, and declared mixed-precision tolerances;
  \item masking, variable lengths, random chunking, and state carry;
  \item cache ties/replacement, segment sealing, live eviction, archive merge,
  and ring wrap;
  \item numerical stress without silent NaN, Inf, or denominator failure.
\end{itemize}

Only after this gate should kernel profiling drive fusion: recurrence scan,
production sliding-window attention, indexed cache insertion/read, packed
episodic reads and routing, and finally compression scheduling.

\subsection{Memory traffic and energy}

DRAM bytes/token must come from hardware counters rather than shape
multiplication. Reports should include DRAM reads/writes, cache hit rates,
achieved bandwidth, occupancy, matrix-unit utilization, launch counts, and the
time share of FFT, QR, and SVD operations.

For sampled device power \(P(t)\), idle power \(P_{\mathrm{idle}}\), and \(N\)
generated tokens, define
\begin{equation}
  E_{\mathrm{token}}=
  \frac{1}{N}\int_{t_0}^{t_1}
  \max\!\left(P(t)-P_{\mathrm{idle}},0\right)\,dt.
  \label{eq:energy}
\end{equation}
The measurement must report sampling interval, integration rule, idle
procedure, power/clock policy, workload duration, and uncertainty. Short
workloads are repeated long enough to exceed sensor resolution.

\subsection{Result-artifact schema}

Every future result artifact should contain:

\begin{itemize}
  \item a versioned protocol ID and source commit or archive hash;
  \item full model, data, optimizer, schedule, precision, and seed
  configuration;
  \item raw per-seed and per-repeat samples;
  \item software, hardware, driver/runtime, thread count, and kernel mode;
  \item parameter count, logical state bytes, allocated peak, and saturation
  test;
  \item explicit classification as smoke, diagnostic, model-only, or
  end-to-end evidence.
\end{itemize}

The shipped CLIs write JSON atomically, reducing the chance that an interrupted
process leaves a partial result that looks complete.
